\subsection{全纯函数演算的存在性与唯一性}

假设 A 是一个复交换幺 Banach 代数。

定义 2. --- 设 \(n \geqslant 1\) 是一个整数，\(\boldsymbol{a} \in \mathrm{A}^{n}\)。设 U 是 \(\mathrm{Sp}^{n}(\boldsymbol{a})\) 的一个开邻域。我们称族 \(\boldsymbol{a}'\) 是 \((\boldsymbol{a}, \mathrm{U})\) 的一个包络，如果 \(\boldsymbol{a}' \in \mathbf{C}^{n+p}\) 延拓 \(\boldsymbol{a}\)，并且 \(\mathrm{U} \times \mathbf{C}^{p}\) 包含 \(\mathrm{Sp}^{n+p}(\boldsymbol{a}')\) 的多项式凸包络。

引理 11. --- 设 \(n \geqslant 1\) 是一个整数。设 \(\boldsymbol{a} \in \mathrm{A}^{n}\)。对 \(\mathrm{Sp}^{n}(\boldsymbol{a})\) 的每个开邻域 U，存在 \((\boldsymbol{a}, \mathrm{U})\) 的一个包络。

设 \((a_{\lambda})_{\lambda\in\Lambda}\) 是 A 的元素所成的一个族，它延拓族 a 并且拓扑生成幺 Banach 代数 A。设 \(\pi\) 是从 \(C^{\Lambda}\) 到 \(C^{n}\) 上的典范投影，并设 \(U'=\pi^{-1}(U)\)。则 \(U'\) 是 \(\mathrm{Sp}^{\Lambda}((a_{\lambda}))\) 的一个邻域，且 \(\mathrm{Sp}^{\Lambda}((a_{\lambda}))\) 是多项式凸的 (I, p.~44, 引理 4)。由 I, p.~47, 引理 6，存在一个有限子集 \(\Lambda_{0}\)（含于 \(\Lambda\) 且包含 \(\{1,2,\ldots,n\}\)），使得 \(\mathrm{pr}_{\Lambda_{0}}(\mathrm{U}')\) 包含 \(\mathrm{pr}_{\Lambda_{0}}(\mathrm{Sp}^{\Lambda}((a_{\lambda})_{\lambda\in\Lambda}))=\mathrm{Sp}^{\Lambda_{0}}((a_{\lambda})_{\lambda\in\Lambda_{0}})\) 的多项式凸包 S。设 \(p\geqslant0\) 是使 \(\Lambda_{0}\) 的基数为 \(n+p\) 的整数，并设 j 是从 \(\{1,\ldots,n+p\}\) 到 \(\Lambda_{0}\) 上的、在 \(\{1,\ldots,n\}\) 上与恒等映射重合的双射。由于 S 的投影包含于 U，族 \((a_{j(k)})_{1\leqslant k\leqslant n+p}\) 是 \((\boldsymbol{a},\mathrm{U})\) 的一个包络。

命题 5. --- 诸映射 \(\Theta_{\boldsymbol{a}}\)（对 \(n \geqslant 1\) 及 \(\boldsymbol{a} \in \mathrm{A}^{n}\)）的给定，构成 A 上的一个全纯函数演算，也就是说，I, p.~51 的条件 (CF1)、(CF2) 与 (CF3) 均得到满足。

设 \(n \geqslant 1\) 是一个整数，\(\boldsymbol{a} = (a_{1}, \ldots, a_{n}) \in \mathrm{A}^{n}\)。映射 \(\Theta_{\boldsymbol{a}}\) 满足 \(\Theta_{\boldsymbol{a}}(z_{i}) = a_{i}\)（对每个满足 \(1 \leqslant i \leqslant n\) 的 i；由 I, p.~64, 引理 10），这就证明了性质 (CF2)。I, p.~63, 引理 9 蕴涵诸映射 \(\Theta_{\boldsymbol{a}}\) 的性质 (CF3)。

映射 \(\Theta_{\pmb{a}}\) 是 A-线性的且连续的 (I, p.~61, 第 5 小节)。它满足 \(\Theta_{\pmb{a}}(1) = 1\) (I, p.~63, 引理 9)。为验证条件 (CF1)，只需再证明 \(\Theta_{\pmb{a}}\) 是一个代数态射。为此，我们将证明：\(\Theta_{\pmb{a}}^{\mathrm{U}}\)（对 \(\mathrm{Sp}^{n}(\pmb{a})\) 的每个开邻域 U）都是代数态射。

先设 U 包含 \(\mathrm{Sp}^{n}(\boldsymbol{a})\) 的多项式凸包 K。设 \(f_{1}\) 与 \(f_{2}\) 是 \(\mathcal{O}(\mathrm{U};\mathrm{A})\) 的元素。存在多项式函数的序列 \((f_{1,k})\)（相应地，\((f_{2,k})\)），它收敛于 \(f_{1}\)（相应地，收敛于 \(f_{2}\)）——收敛是在 \(\mathcal{O}(\mathrm{K};\mathrm{A})\) 中进行的 (I, p.~68, 定理 2)，因而在 \(\mathcal{O}(\mathrm{Sp}^{n}(\boldsymbol{a});\mathrm{A})\) 中收敛。对每个整数 \(k\)，有

\[
\Theta_ {\boldsymbol {a}} ^ {\mathrm{U}} (f _ {1, k}) \Theta_ {\boldsymbol {a}} ^ {\mathrm{U}} (f _ {2, k}) = \Theta_ {\boldsymbol {a}} ^ {\mathrm{U}} (f _ {1, k} f _ {2, k})
\]

（由 I, p.~64, 引理 10），取极限便得 \(\Theta_{a}^{\mathrm{U}}(f_{1})\Theta_{a}^{\mathrm{U}}(f_{2}) = \Theta_{a}^{\mathrm{U}}(f_{1}f_{2})\)。

考察一般情形。设 \(a' \in C^{n+p}\) 是 \((\boldsymbol{a}, U)\) 的一个包络（引理 11），\(\pi: U \times C^{p} \to U\) 是典范投影。由于 \(U \times C^{p}\) 包含 \(\mathrm{Sp}^{n+p}(\boldsymbol{a}')\) 的多项式凸包络，有

\[
\Theta_ {\boldsymbol {a} ^ {\prime}} ^ {\mathrm{U} \times \mathbf {C} ^ {p}} (f _ {1} \circ \pi) \Theta_ {\boldsymbol {a} ^ {\prime}} ^ {\mathrm{U} \times \mathbf {C} ^ {p}} (f _ {2} \circ \pi) = \Theta_ {\boldsymbol {a} ^ {\prime}} ^ {\mathrm{U} \times \mathbf {C} ^ {p}} (f _ {1} f _ {2} \circ \pi)
\]

（对 \(f_1\) 与 \(f_2\)——均属于 \(\mathcal{O}(\mathrm{U};\mathrm{A})\)——由第一种情形得到）。又因对每个函数 \(f \in \mathcal{O}(\mathrm{U};\mathrm{A})\)，有 \(\Theta_{\boldsymbol{a}'}^{\mathrm{U}\times \mathbf{C}^p}(f\circ \pi) = \Theta_{\boldsymbol{a}}^{\mathrm{U}}(f)\)（前面已证的条件 (CF3)），结论成立，从而条件 (CF1) 成立。

现在我们可以证明 I, p.~51 的定理 1。命题 5 表明映射族 \((\Theta_{a})_{a}\) 是 A 上的一个全纯函数演算。因此只需再确立 A 上全纯函数演算的唯一性。

设 \((\Psi_{a})_{a}\) 是一个映射族，它对每个整数 \(n \geqslant 1\) 以及每个 \(a \in A^{n}\) 有定义，并满足 A 上全纯函数演算的条件 (CF1)、(CF2)、(CF3) (I, p.~51)。只需证明：对每个整数 \(n \geqslant 1\)、每个 \(a \in A^{n}\) 以及 a 的每个开邻域 U，我们有 \(\Theta_{a}^{U} = \Psi_{a}^{U}\)。

设 \(n \geqslant 1\)，\(\boldsymbol{a} = (a_{1}, \ldots, a_{n}) \in \mathrm{A}^{n}\)。设 U 是 \(\mathrm{Sp}^{n}(\boldsymbol{a})\) 的一个开邻域。先设 U 包含 \(\mathrm{Sp}^{n}(\boldsymbol{a})\) 的多项式凸包 K。由性质 (CF1) 与 (CF2)，态射 \(\Theta_{a}^{U}\) 与 \(\Psi_{a}^{U}\) 在多项式函数上相一致。由 I, p.~68, 定理 2 的推论以及连续性性质 (CF1)，这两个态射因此相等。

我们来证明一般情形。设 \(a' \in C^{n+p}\) 是 \((\boldsymbol{a}, U)\) 的一个包络，\(\pi: U \times C^{p} \to U\) 是典范投影。我们有

\[
\Theta_ {\boldsymbol {a}} ^ {\mathrm{U}} = \Theta_ {\boldsymbol {a} ^ {\prime}} ^ {\mathrm{U} \times \mathbf {C} ^ {p}} \circ \pi^ {*} = \Psi_ {\boldsymbol {a} ^ {\prime}} ^ {\mathrm{U} \times \mathbf {C} ^ {p}} \circ \pi^ {*} = \Psi_ {\boldsymbol {a}} ^ {\mathrm{U}}
\]

（由性质 (CF3) 及前面的情形）。这就完成了 I, p.~51 定理 1 的证明。

注意，I, p.~68 的定理 2 还蕴涵下述唯一性结果：

命题 6. --- 设 \(\boldsymbol{a} \in \mathrm{A}^n\)。我们假设 \(\mathrm{Sp}^n(\boldsymbol{a})\) 是多项式凸的。设 \(z_1, \ldots, z_n\) 是 \(\mathrm{Sp}^n(\boldsymbol{a})\) 的某个邻域中的、\(\mathbf{C}^n\) 上坐标函数的芽。则映射 \(\Theta_{\boldsymbol{a}}\) 是唯一的连续幺代数态射 \(\varphi\)：从 \(\mathcal{O}(\mathrm{Sp}^n(\boldsymbol{a}); \mathrm{A})\) 到 A 中，使得 \(\varphi(z_1) = a_1, \ldots, \varphi(z_n) = a_n\)。

I, p.~64 的引理 10 与 I, p.~65 的命题 2 的推论为全纯函数演算论证了下述记号的合理性。设 \(n \geqslant 1\)，\(\boldsymbol{a} \in \mathrm{A}^n\)。对每个芽 \(f \in \mathcal{O}(\mathrm{K};\mathrm{A})\)（相应地，对每个全纯函数 \(f \in \mathcal{O}(\mathrm{U};\mathrm{A})\)——它定义在 \(\mathrm{Sp}^n (\boldsymbol{a})\) 的某个开邻域 U 上），我们令

\[
f (\boldsymbol {a}) = \Theta_ {\boldsymbol {a}} (f).\tag{6}
\]

依性质 (CF1) 与 (CF2)，当 f 是多项式时，这一记号与 A, IV, p.~4, n°3 中引入的记号相容。于是 I, p.~51 的性质 (CF2) 与 (CF3) 可以写成

\[
z _ {i} (\boldsymbol {a}) = a _ {i}, \quad 1 \leqslant i \leqslant n, \quad (f \circ \pi_ {m, n}) (\boldsymbol {a}) = f (\pi_ {m, n} (\boldsymbol {a})).
\]

